\section{Introduction}

AI-Scientist belongs to a broader family of research systems that seek to improve the reliability of language-model-based reasoning. The related literature spans several areas: hallucination in large language models, retrieval-augmented generation, scientific claim verification, self-critique, and multi-agent orchestration \cite{ji2023hallucination,huang2023hallucinationsurvey,ragsurvey2024}. For a thesis of this kind, the literature review benefits from doing more than listing paper titles. It explains how these lines of work fit together, what they already demonstrate, and why a verification-first framework offers a meaningful contribution.

The central idea emerging from the literature is encouraging: strong research assistants do not need to be perfect generators, but they do need to be disciplined evidence users. This thesis adopts that perspective and positions AI-Scientist as a system that is designed to retrieve evidence, isolate claims, verify those claims, critique weak conclusions, and present the final result in a transparent way.

\begin{figure}[H]
\centering
\begin{tikzpicture}[
    scale=0.96,
    transform shape,
    node distance=1.2cm and 1.0cm,
    box/.style={rectangle, rounded corners, draw=black, align=center, minimum width=3.2cm, minimum height=0.95cm, fill=blue!18},
    claim/.style={rectangle, rounded corners, draw=black, align=center, minimum width=3.7cm, minimum height=0.95cm, fill=gray!30},
    synth/.style={rectangle, rounded corners, draw=black, align=center, minimum width=3.4cm, minimum height=1.0cm, fill=yellow!25},
    arrow/.style={-{Latex[length=3mm]}, thick}
]
\node[box] (a) {Hallucination\\risk};
\node[box, right=of a] (b) {Grounded\\retrieval};
\node[claim, right=of b] (c) {Claim-level\\verification};
\node[synth, below=of b] (d) {Multi-agent\\critique and synthesis};
\draw[arrow] (a) -- (b);
\draw[arrow] (b) -- (c);
\draw[arrow] (c) -- (d);
\draw[arrow] (b) -- (d);
\end{tikzpicture}
\caption{Literature-to-design mapping for AI-Scientist}
\label{fig:lit_to_design}
\end{figure}

\section{Hallucination, Grounding, and Retrieval}

The first major theme in the literature is hallucination. In the context of scientific question answering, hallucination refers to fluent but unsupported output: a model may present a claim confidently even when the evidence does not justify it. This is especially problematic in research settings because incorrect claims can distort interpretation, misrepresent prior work, or create the impression of certainty where the literature is actually mixed \cite{manakul2023selfcheckgpt,min2023factscore,ji2023hallucination}. The literature therefore strongly supports the idea that trustworthy systems must expose their evidence rather than hiding it.

Retrieval-augmented generation addresses this issue by connecting a model to external text \cite{lewis2020rag,guu2020realm}. Instead of relying only on internal memory, the system fetches documents or passages that are relevant to the question and conditions the answer on those passages. This improves groundedness and gives the model a stronger factual basis. The key lesson for AI-Scientist is that retrieval is necessary, while effective verification still depends on interpretation, comparison, and support checking.

That observation is important because scientific tasks are rarely solved by quoting one passage. A real answer may depend on multiple papers, overlapping claims, and subtle differences in experimental setup. For that reason, the thesis treats retrieval as the first step in a wider verification pipeline rather than as the final answer mechanism. This framing keeps the system research-oriented and makes the later verification and critique stages more meaningful.

\begin{table}[H]
\centering
\begin{tabularx}{\linewidth}{>{\raggedright\arraybackslash}p{0.28\linewidth} >{\raggedright\arraybackslash}X}
\toprule
System idea & What it contributes to this thesis \\
\midrule
Plain generation & Produces fluent answers with limited evidence traceability \\
RAG & Adds external evidence and improves grounding \\
Retrieval + ranking & Helps choose the most relevant source passages \\
AI-Scientist & Uses retrieval as the input to claim verification, critique, and synthesis \\
\bottomrule
\end{tabularx}
\caption{How retrieval-oriented systems differ in the thesis context}
\end{table}

\section{Claim Verification and Structured Evaluation}

The second major theme is scientific claim verification \cite{thorne2018fever,wadden2020scifact,honovich2022true}. This literature is especially important because it matches the unit of analysis used by AI-Scientist. Instead of evaluating an entire answer as a single blob of text, claim verification decomposes the answer into smaller statements and checks each one against evidence. This is a stronger and more research-friendly approach because support, contradiction, and uncertainty are often mixed together in a single response.

Claim-level verification also makes evaluation easier to interpret. If one claim is supported and another is contradicted, the system can report that difference explicitly rather than hiding it in an averaged score. That is why this thesis later uses not only overall accuracy, but also claim-level precision, recall, and F1-score. These metrics are particularly useful because they capture how well the system identifies supported claims, how often it misses them, and how cleanly it separates correct from incorrect verdicts.

Another useful insight from the literature is that evidence should not only be present; it should be attributable. A verdict without evidence is not very useful in a scientific setting because the reader cannot inspect why the system chose that conclusion. The fact-checking and claim-verification literature consistently emphasizes rationales, supporting sentences, and explicit labels such as supported, contradicted, or insufficient evidence \cite{wadden2020scifact,thorne2018fever,min2023factscore}. AI-Scientist adopts the same idea and applies it in an orchestrated setting.

\section{Self-Critique, Multi-Agent Reasoning, and Orchestration}

The third major theme is self-critique and multi-step reasoning. Single-pass generation is fast, but it often benefits from an additional layer of review when research-grade support is required. The literature shows that systems improve when they are allowed to revise their own outputs, ask follow-up questions, or re-examine weak statements before finalizing an answer \cite{shinn2023reflexion,asai2024selfrag,yao2023treeofthoughts}. This thesis follows that idea in a structured way, using a dedicated critic stage to make the reasoning process more disciplined and inspectable.

Multi-agent orchestration is the natural way to organize that process \cite{wu2023autogen,li2023camel,du2023debate}. A retrieval agent can focus on finding evidence, a claim-extraction agent can break the evidence into checkable statements, a verification agent can decide whether the statements are supported, and a critic agent can re-check weak or contradictory claims. Finally, a synthesis agent can assemble a coherent response. This separation of responsibilities is valuable because it makes the reasoning process clearer and the system behavior easier to analyze.

The important point from the literature is that multiple agents become especially valuable when their roles are specialized. When each agent has a focused purpose, the full pipeline becomes more disciplined and more inspectable. That makes the thesis stronger because the architecture is presented as a deliberate verification workflow with clear research value.

\begin{figure}[H]
\centering
\begin{tikzpicture}[
    scale=0.96,
    transform shape,
    node distance=1.1cm,
    box/.style={rectangle, rounded corners, draw=black, align=center, minimum width=3.3cm, minimum height=0.92cm, fill=blue!18},
    critique/.style={rectangle, rounded corners, draw=black, align=center, minimum width=3.3cm, minimum height=0.92cm, fill=yellow!25},
    synth/.style={rectangle, rounded corners, draw=black, align=center, minimum width=3.3cm, minimum height=0.92cm, fill=orange!22},
    arrow/.style={-{Latex[length=3mm]}, thick}
]
\node[box] (r) {Retrieve evidence};
\node[box, below=of r] (e) {Extract claims};
\node[box, below=of e] (v) {Verify claims};
\node[critique, below=of v] (c) {Critique weak claims};
\node[synth, below=of c] (s) {Synthesize final report};
\draw[arrow] (r) -- (e);
\draw[arrow] (e) -- (v);
\draw[arrow] (v) -- (c);
\draw[arrow] (c) -- (s);
\end{tikzpicture}
\caption{Conceptual multi-agent workflow reflected in AI-Scientist}
\label{fig:multi_agent_workflow}
\end{figure}

\section{Implications for the Thesis Design}

The literature reviewed above leads to several practical implications for AI-Scientist. First, grounding should be visible, not hidden. Second, verification should happen at the claim level, not only at the answer level. Third, critique should be a deliberate stage in the pipeline so that weak conclusions can be revisited before final reporting. Fourth, confidence should be explicit because it helps the user understand which conclusions are strongly supported and which are merely plausible \cite{kadavath2022language,lin2021truthfulqa}. Finally, the complete pipeline should be traceable so that the reasoning can be explained in a thesis, defended in a viva, and inspected by another researcher.

These implications shape the structure of the remaining chapters. The architecture chapter will show how the agents are connected. The implementation chapter will explain how each agent behaves. The evaluation chapter will then test whether the framework produces strong results on the benchmark and whether the critic loop improves the outcome.

\section{Comparison of Themes}

Table~\ref{tab:related_work_comparison} provides a compact comparison of the literature themes and their relevance to the thesis.

\begin{table}[H]
\centering
\begin{tabularx}{\linewidth}{>{\raggedright\arraybackslash}p{0.22\linewidth} >{\raggedright\arraybackslash}p{0.23\linewidth} >{\raggedright\arraybackslash}X}
\toprule
Theme & Main strength & Thesis relevance \\
\midrule
Hallucination studies & Identify failure modes in model output & Justify the need for evidence grounding \\
RAG & Bring external evidence into generation & Provide the retrieval backbone \\
Claim verification & Judge claims against evidence & Enable claim-level verdicts and metrics \\
Self-critique & Revise weak outputs through feedback & Support re-verification of uncertain claims \\
Multi-agent orchestration & Separate responsibilities across stages & Make the full workflow transparent \\
\bottomrule
\end{tabularx}
\caption{Comparison of related-work themes}
\label{tab:related_work_comparison}
\end{table}

\section{Chapter Summary}

This chapter showed that AI-Scientist is grounded in a strong and coherent research tradition. The literature supports the need for grounded retrieval, claim-level verification, explicit confidence, and structured critique. It also shows why a multi-agent architecture is a natural fit for the problem: the system needs to separate evidence finding, claim checking, critique, and synthesis into distinct stages. The next chapter formalizes the problem statement and explains the requirements of the proposed framework in more detail.
